\documentclass{beamer}
\usetheme{Madrid}
\usecolortheme{default}

\title[Spiking Neural Networks]{\textbf{SPIKING NEURAL NETWORKS FOR MULTIMODAL NEUROIMAGING}}
\subtitle{An Analysis of Current Trends and the NeuCube Brain-Inspired Architecture}
\author{
\textbf{Submitted By:} \\ SUHAIL K B (CEC22CS123) \\[0.2cm]
\textbf{Guide Name:} \\ Mrs. ANSILA A S \\ Assistant Professor
}
\institute{
Department of Computer Engineering \\ College of Engineering, Cherthala}
\date{October 21, 2025}

\begin{document}

%-------------------- TITLE SLIDE --------------------
\begin{frame}
  \titlepage
\end{frame}

%-------------------- INTRODUCTION --------------------
\begin{frame}{Introduction}
\begin{itemize}
  \item AI has revolutionized neuroimaging for disease detection and multimodal data integration.
  \item ANNs face challenges with temporal dynamics and computational cost.
  \item SNNs mimic biological neurons using discrete spikes.
  \item NeuCube is a specialized SNN framework for spatiotemporal brain data.
\end{itemize}
\end{frame}

%-------------------- LITERATURE SURVEY --------------------
\begin{frame}{Literature Survey}
\begin{enumerate}
  \item \textbf{Improving NeuCube (2023):} Enhanced EEG pattern recognition via transfer learning.
  \item \textbf{ANN-SNN Conversion (2023):} Retains ANN accuracy with SNN efficiency.
  \item \textbf{BackEISNN (2022):} Balanced excitatory-inhibitory neurons for adaptive learning.
  \item \textbf{SpikingJelly (2023):} Python-based framework for spike-based intelligence.
\end{enumerate}
\end{frame}

%-------------------- FUNDAMENTALS --------------------
\begin{frame}{Fundamentals of SNNs and Neuromorphic Computing}
\textbf{Neuromorphic Computing:}Brain-inspired computing: High energy efficiency via parallel, event-driven processing.\\
\medskip
\textbf{Spiking Neural Networks (SNNs):} Use discrete, neuron-like spikes to process dynamic, time-varying data.\\
\medskip
\textbf{Key Concepts:}
\begin{itemize}
    \item \textbf{Spikes:} 
    \item \textbf{Membrane Potential:} 
    \item \textbf{Temporal Coding:} 
    \item \textbf{Synaptic Plasticity:} 
    \item \textbf{Refractory Period:} 

\end{itemize}
\end{frame}
%-------------------- SNN vs ANN --------------------
\begin{frame}{SNN Neuron Models}
    \begin{itemize}
        \item \textbf{Leaky Integrate-and-Fire (LIF):}
            \begin{itemize}
                \item \textbf{Integrate:} Input spikes increase membrane potential.
                \medskip
                \item \textbf{Leak:} Adds biological realism; potential decays over time. 
                \medskip
                \item \textbf{Fire:} Emits a spike when potential reaches threshold.
                \medskip
                \item \textbf{Reset:} Potential resets after firing.
            \end{itemize}
    \end{itemize}
\end{frame}
%-------------------- NEUCUBE --------------------
\begin{frame}{SNN Neuron Models}

\begin{figure}[h]
  \centering
  \includegraphics[width=0.8\linewidth]{fig3.jpg} % replace with your image filename
\end{figure}
\end{frame}


%-------------------- SNN vs ANN --------------------
\begin{frame}{SNN Neuron Models}
    \begin{itemize}
        \item \textbf{Integrate-and-Fire (IF):} The simplest model that integrates incoming current.
        \medskip
        \item \textbf{Izhikevich:} Balances biological accuracy with computational efficiency.
        \medskip
        \item \textbf{Hodgkin-Huxley:} The most detailed and accurate model, but computationally heavy.
        \medskip
        \item \textbf{AdEx:} An advanced model that mimics neuron fatigue.
    \end{itemize}
\end{frame}
    


%-------------------- SNN vs ANN --------------------
\begin{frame}{SNN Learning Models}
    \begin{itemize}
        \item \textbf{Unsupervised :} based on the relative timing of pre- and post-synaptic spikes.
        \medskip
        \begin{enumerate}
            \item \textbf{Spike-Timing Dependent Plasticity (STDP)}\\
            If a presynaptic neuron fires \textbf{before} a postsynaptic neuron → \textbf{strengthen} the synapse (Long-Term Potentiation).\\
            \medskip
            If it fires \textbf{after} → \textbf{weaken} the synapse (Long-Term Depression).\\
        \medskip
            \item \textbf{Hebbian Learning}

            “Neurons that fire together, wire together.”\\
        \end{enumerate}

    \end{itemize}
\end{frame}

%-------------------- SNN vs ANN --------------------
\begin{frame}{SNN Learning Models}
    \begin{itemize}
        \item \textbf{Supervised Learning using Surrogate Gradients}        \begin{itemize}
            \item Surrogate derivatives are used to approximate gradients.
            \item Enables backpropagation-like training for SNNs.
            \item Trained on labeled datasets like ANNs.
        \end{itemize}
        \bigskip
        \item \textbf{Reinforcement (R-STDP):} 
        \begin{itemize}
            \item Synaptic weights based on reward signals.
            \item Combines spike activity with feedback from the environment.
            \item Useful for robotics or neuromorphic control tasks.
        \end{itemize}
    \end{itemize}
\end{frame}

%-------------------- SNN vs ANN --------------------
\begin{frame}{Spike Encoding}
    SNNs encode information in spikes, mainly in two ways:
    \begin{itemize}
        \item \textbf{Rate coding:} Number of spikes in a time window.\\
        Example: More spikes → higher intensity.
        \medskip
        \item \textbf{Temporal coding:} Exact timing of spikes.\\
        Example: The first neuron to fire represents the strongest input.
    \end{itemize}

\end{frame}

%-------------------- SNN vs ANN --------------------
\begin{frame}{Spike Encoding}
\begin{figure}[h]
\centering
  \includegraphics[width=0.8\linewidth]{fig4.jpg} % replace with your image filename
\end{figure}
\begin{itemize}
    \item Integrate incoming spikes → accumulate membrane potential.
    \item Fire when threshold is reached → send spikes.
    \item Learning adjusts weights based on spike timing or approximated gradients
\end{itemize}
\end{frame}

%-------------------- SNN vs ANN --------------------
\begin{frame}{SNNs vs ANNs}
\begin{figure}[h]
  \centering
  \includegraphics[width=1\linewidth]{image.png} % replace with your image filename
\end{figure}
\end{frame}

%-------------------- APPLICATIONS --------------------
\begin{frame}{Applications in Neuroimaging}
\begin{itemize}
    \item \textbf{Structural Neuroimaging (MRI):} 
    \medskip
    \item \textbf{Functional Neuroimaging (EEG/fMRI):} 
    \textbf{Applications:} Brain-Computer Interfaces (BCIs), disease diagnosis (e.g., Alzheimer's, epilepsy), and classifying cognitive states.
    \medskip
    \item \textbf{Multimodal Data Integration:}  
\end{itemize}
\end{frame}

%-------------------- DEVELOPMENT ECOSYSTEM --------------------
\begin{frame}{SNN Development Ecosystem}
\par \textbf{Development Challenges:} Adoption is hindered by a lack of pre-built models, specialized software, and standardized formats.

\textbf{Software Simulators:}
\begin{itemize}
  \item snnTorch
  \item SpikingJelly
  \item Nengo DL
  \item NeuCube
\end{itemize}

\end{frame}

%-------------------- NEUCUBE --------------------
\begin{frame}{Industry Hardwares}
\textbf{Hardware Platforms:}
\begin{itemize}
  \item \textbf{Intel Loihi2 :}A research chip with programmable neuron models and on-chip learning. 1.15 billion neuron system based on this architecture;95\% accuracy in gesture recognition with 200x lower power than a GPU.
   \medskip
  \item \textbf{IBM NorthPole:} A chip focused on high-performance, energy-efficient model inference, rather than SNN development.
  \medskip
  \item \textbf{BrainChip Akida: } Paired with a CPU for hybrid execution.includes a CNN-to-SNN conversion tool. The platform has shown 98\% accuracy recognizing fast-moving objects on edge devices.
  \medskip
  \item \textbf{Lynxi KA200: } A high-performance AI accelerator designed for edge computing and autonomous driving.
\end{itemize}
\end{frame}


%-------------------- NEUCUBE --------------------
\begin{frame}{NeuCube: Architecture \& Implementation}
\begin{itemize}
  \item Brain-inspired 3D structure for spatiotemporal mapping.
  \item Integrates multimodal data.
  \item Supports hybrid learning (unsupervised + supervised).
\end{itemize}
\textbf{Architectural Components:}
\begin{itemize}
    \item \textbf{Input Data Stream:} 
    \item \textbf{Gene Regulatory Network (GRN):}
    \item \textbf{Neurogenetic Cube:} 
    \item \textbf{Output Module:} 
\end{itemize}
\end{frame}

%-------------------- NEUCUBE --------------------
\begin{frame}{The NeuCube Framework}

\begin{figure}[h]
  \centering
  \includegraphics[width=1\linewidth]{fig2.jpg} % replace with your image filename
\end{figure}
\end{frame}

\begin{frame}{The NeuCube Framework}

\textbf{Implementation \& Flexibility:}
\begin{itemize}
    \item NeuCube supports flexible implementation on conventional computers, GPUs, and neuromorphic hardware (e.g., Loihi and SpiNNaker).
    \bigskip
    \item Integrates diverse data, adapt its structure, and incorporate GRNs makes it a versatile hybrid platform.
\end{itemize}
\end{frame}



%-------------------- CASE STUDIES --------------------
\begin{frame}{Case Study 1: MRI Classification}
\begin{itemize}
  \item Compared SNN (SpikingJelly) and CNN (AlexNet) for Alzheimer’s detection.
  \item CNN: 85.9\% accuracy vs SNN: 78\% (faster but less precise).
  \item SNNs less effective for static MRI data.
\end{itemize}


\begin{figure}[h]
  \centering
  \includegraphics[width=0.7\linewidth]{fig 1.jpg} % replace with your image filename
  \caption{Accuracy comparison of CNN and SNN on MRI data}
\end{figure}

\end{frame}






\begin{frame}{Case Study 2: EEG and fMRI Analysis}
\begin{itemize}
  \item NeuCube achieved 97\% accuracy in EEG-based disease detection.
  \item fMRI: 90\% accuracy in cognitive state classification.
  \item Outperformed traditional ML models for temporal data.
\end{itemize}
\end{frame}

\begin{frame}{Case Study 3: Longitudinal MRI Data}
\begin{itemize}
  \item NeuCube trained on 6-year MRI dataset for dementia prediction.
  \item Predicted disease 2 years in advance with 91\% accuracy.
  \item Demonstrates potential for long-term predictive neuroimaging.
\end{itemize}
\end{frame}

%-------------------- HYBRID SYSTEMS --------------------
\begin{frame}{Hybrid Systems: Integrating ANNs and SNNs}
\begin{itemize}
  \item Combines CNN’s spatial analysis with SNN’s temporal efficiency.
  \item Used in neuroprosthetics and adaptive neuro-fuzzy systems.
  \item Enhances brain-computer interface performance.
\end{itemize}
\end{frame}

%-------------------- CHALLENGES --------------------
\begin{frame}{Challenges and Future Outlook}
\begin{itemize}
  \item Lack of standardization and benchmarking.
  \item Sensitive to noise in EEG data.
  \item Difficult training due to non-differentiable spikes.
  \item Ethical issues: privacy, bias, and clinical safety.
  \item \textbf{Future:} predictive modeling, personalized medicine, and BCIs.\\
  \medskip
  Research will focus on predictive modeling for early diagnosis, personalized medicine, and advanced Brain-Computer Interfaces.

\end{itemize}
\end{frame}

%-------------------- CONCLUSION --------------------
\begin{frame}{Conclusion}
\begin{itemize}
  \item SNNs outperform ANNs in temporal neuroimaging tasks like EEG/fMRI.
  \item NeuCube demonstrates strong multimodal integration and adaptability.
  \item Hybrid ANN-SNN models show real-world potential.
  \item Standardization and clinical validation are key future steps.
\end{itemize}
\end{frame}

%-------------------- REFERENCES --------------------
\begin{frame}{References}
\scriptsize
\begin{thebibliography}{6}
\bibitem[1]{kasabov2014} N. K. Kasabov, ``NeuCube: A spiking neural network architecture for mapping and learning spatio-temporal brain data,'' \textit{Neural Networks}, 2014.
\bibitem[2]{kasabov2015} N. Kasabov \& E. Capecci, ``Spiking neural network methodology for EEG data modelling,'' \textit{Information Sciences}, 2015.
\bibitem[3]{doborjeh2021} M. Doborjeh et al., ``Personalised Predictive Modelling with SNNs of MRI Data,'' \textit{Neural Networks}, 2021.
\bibitem[4]{nunes2022} J. Nunes et al., ``Spiking Neural Networks: A Survey,'' \textit{IEEE Access}, 2022.
\bibitem[5]{wu2023} X. Wu et al., ``Improving NeuCube for EEG Pattern Recognition,'' \textit{Neurocomputing}, 2023.
\bibitem[6]{pfeiffer2018} M. Pfeiffer \& T. Pfeil, ``Deep Learning with Spiking Neurons: Opportunities and Challenges,'' \textit{Frontiers in Neuroscience}, 2018.
\end{thebibliography}
\end{frame}

\end{document}